%% ============================================================
%% Chapter 18: Clifford FHE as the Inverse Compression Case
%% ============================================================

\chapter{Clifford FHE as the Inverse Compression Case}
\label{ch:ch18-clifford-fhe}
\section{Introduction}

Chapter~\ref{ch:ch15-compression-fallacy} defined a representation's three coordinates, description magnitude $L_T$, distinguishing structure $\mathcal{P}_T$, and operational repertoire $\mathcal{O}_T$, and showed the first two are independent while the third is a genuinely separate axis a lossless archive can leave essentially empty. This chapter closes Part~IV with the sharpest available instance of that third coordinate: a case, drawn from an external, independently authored paper on homomorphic encryption, in which two encodings agree completely on stored values and on static distinguishing capacity while disagreeing entirely on what can be computed over those values without first decrypting them. The reading below is offered in the same spirit as Chapter~\ref{ch:ch08-tartan-tiling}'s treatment of Buchanan's contextuality paper: a structural correspondence recognized after the fact, not a claim that the source paper was written with this monograph's apparatus in mind, and every quantitative claim below is attributed to its source rather than adopted as this monograph's own.

\section{The Problem: Flattening Discards Structure}

Conventional fully homomorphic encryption (FHE) schemes support arbitrary computation on encrypted data without decryption, but the standard CKKS scheme, which supports approximate arithmetic on real numbers, is scalar-only: it has no native awareness of geometric relationships and requires flattening a structured object, such as an eight-component multi-vector, into eight separate, uncoordinated ciphertexts. Geometric algebra represents geometric objects through multi-vectors, elements combining scalars, vectors, bivectors, and higher grades within one unified algebraic framework. In the three-dimensional Euclidean geometric algebra $\mathrm{Cl}(3,0)$, a general multi-vector
\[
M = m_0 + m_1 e_1 + m_2 e_2 + m_3 e_3 + m_{12}e_{12} + m_{13}e_{13} + m_{23}e_{23} + m_{123}e_{123}
\]
combines a scalar, three vector components, three bivector components representing oriented planes, and one trivector component representing oriented volume, eight components in total. The defining operation is the geometric product $ab = a\cdot b + a\wedge b$, associative, distributive, and generally non-commutative, from which rotations are realized via rotor conjugation, $v' = RvR^{-1}$ (written $R \otimes v \otimes \tilde R$ in the source paper's notation), composing more compactly and more cheaply than matrix multiplication.

The source paper's central question is whether geometric data can be encrypted while preserving its geometric properties, rather than being flattened into scalar arithmetic that discards them. It answers with Clifford FHE, the first ring-learning-with-errors-based FHE scheme with native support for all seven fundamental Clifford algebra operations (geometric product, reverse, rotation, wedge product, inner product, projection, and rejection), and with geometric neural networks that operate directly on encrypted multi-vectors, replacing matrix multiplication with the geometric product as the network's fundamental layer operation.

\section{Two Encodings, Identical Values, Different Repertoires}

It is tempting to read the headline figure the source paper reports, a $250{,}000\times$ expansion between a 64-byte plaintext multi-vector and its roughly 16-megabyte encrypted form, as a direct instance of the coding-deficit gap Chapter~\ref{ch:ch15-compression-fallacy} measures. That reading does not hold up. The expansion is ciphertext expansion, redundancy that ring-learning-with-errors encryption introduces for cryptographic security, not a comparison of description lengths in the sense $\delta_T(x) = L_T(x) - L^\ast(x)$ is defined in this monograph. A large ciphertext is not evidence of a large coding deficit, and the size figures should not be read as one.

The more interesting distinction survives this correction, but it belongs to the third coordinate of Chapter~\ref{ch:ch15-compression-fallacy}'s apparatus, the operational repertoire $\mathcal{O}_T$, not to $L_T$ or $D_T$. Two schemes for encrypting a multi-vector, ordinary flattened CKKS and Clifford FHE, both encrypt each of the eight components separately, eight ciphertexts either way; by the definitions of Chapter~\ref{ch:ch15-compression-fallacy}, $D_T$ is unaffected by this choice, since the component values are equally present and equally distinct under both schemes. What differs is $\mathcal{O}_T$: Clifford FHE additionally encodes a fixed table of structure constants, $C_{ijk} \in \{-1,0,1\}$ specifying how the eight encrypted components combine under the geometric product, so that rotation, projection, and the wedge and inner products become native operations over the ciphertexts, computable directly rather than only after decryption. The homomorphic geometric product computes
\[
\operatorname{Enc}(c_k) = \sum_{i,j} C_{ijk} \cdot \operatorname{Enc}(a_i) \cdot \operatorname{Enc}(b_j),
\]
64 ciphertext multiplications per geometric product, against 512 for the equivalent naive $8\times 8$ matrix encoding, an eightfold reduction in operation count that follows from the geometric product's structural sparsity, only 64 of the 512 possible term products are nonzero, rather than from any change in what is stored.

``Discards'' overstates what ordinary flattened CKKS does to these relations: nothing about scalar CKKS makes a rotation mathematically unrecoverable from the eight ciphertexts, since the component values are all still there and could in principle be recombined by hand using ordinary scalar homomorphic operations. What ordinary CKKS lacks is the Clifford structure as a native, first-class part of $\mathcal{O}_T$: recovering a rotation from flattened ciphertexts requires re-deriving, term by term, an algebra that Clifford FHE simply supplies as an available operation. The source paper's own security proof makes the boundary of what is and is not preserved precise: Clifford FHE's security reduces to CKKS security, since each multi-vector component is an independently encrypted CKKS ciphertext, so an adversary breaking Clifford FHE with advantage $\epsilon$ can be used to break ordinary CKKS with advantage $\epsilon/8$. The two schemes are cryptographically equivalent at the level of what is hidden; they differ only in what is computable without first removing that hiding.

\section{Clifford FHE as the Inverse Compression Case}

Chapter~\ref{ch:ch15-compression-fallacy}'s three coordinates, description magnitude, distinguishing structure, and operational repertoire, define a Pareto-style tradeoff space, and Clifford FHE occupies an extreme and clarifying point in it. Standard CKKS optimizes for description length by flattening a multivector into independent scalar ciphertexts, discarding, or more precisely relocating outside $\mathcal{O}_T$, the algebraic structure of the geometric product. Clifford FHE represents the exact inverse case. It abandons compression entirely, accepting the full $250{,}000\times$ ciphertext expansion the source paper reports, precisely to preserve the algebraic structure of geometric operations natively over the ciphertexts. Thus Clifford FHE is the most mathematically extreme realization of the compression fallacy's central thesis: to preserve native, first-class operational structure, a representation may need to sacrifice essentially all of the compactness that would ordinarily be read as evidence of good design. Representational quality is a genuinely multi-dimensional frontier, $R = L_T \times \mathcal{P}_T \times \mathcal{O}_T$, where size and structure can stand in direct, structural tension rather than trading off along one shared axis.

This also gives Chapter~\ref{ch:ch15-compression-fallacy}'s intervention-robustness apparatus its cleanest available instance. Rotational equivariance in the geometric neural network described below is not an empirical regularity discovered by training on many rotated examples; it follows from the algebraic construction itself. For a rotor $R$ and geometric-linear-layer weight $W$, applying $W \otimes -$ to a rotated input commutes with rotation by associativity of the geometric product and the properties of rotor conjugation:
\[
W \otimes (RxR^{-1}) = R(W \otimes x)R^{-1},
\]
so the layer's output rotates consistently with its input as an algebraic identity, independent of any particular learned weights. In Chapter~\ref{ch:ch15-compression-fallacy}'s notation, $R_T(x,I) = 0$ for $I$ a rotation is a claim provable from $T$'s algebra rather than measured empirically on a held-out test set: a representation's distinguishing capacity, tested under deliberate intervention, is here revealed by proof rather than by observation, exactly the sharper standard that chapter's closing section recommended over trusting a static count taken at rest.

\section{The Geometric Neural Network and Its Reported Results}

The source paper's geometric neural network replaces the ordinary layer $y = Wx + b$ with $y = W \otimes x + b$, where $W$, $x$, and $b$ are all multi-vectors, giving a coordinate-free representation whose rotational equivariance is structural rather than learned. Because homomorphic encryption supports polynomial operations rather than arbitrary nonlinearities such as ReLU, the activation function is normalization, $\mathrm{Act}(m) = m/\|m\|$, approximated homomorphically by a low-degree polynomial. A three-layer network, one input multi-vector through 16 then 8 hidden multi-vectors to a small number of output classes, was evaluated on encrypted three-dimensional point-cloud classification across three shape categories, spheres, cubes, and pyramids, each point cloud reduced to a single multi-vector whose scalar component encodes average radial distance, vector components encode centroid position, bivector components encode principal orientation planes, and trivector component encodes a volume measure.

The reported results are a 99\% encrypted classification accuracy against 100\% on the same task in plaintext, three misclassifications out of 300 encrypted samples, all near class boundaries, with inference completing in under 60 seconds per sample despite the full ciphertext expansion. The reported baseline comparison, a classical multilayer perceptron on three mean-coordinate features reaching 35\% accuracy and failing under rotation, against 99\% for the Clifford encoding, is not a controlled ablation of flattening alone. The two systems receive different input features: the baseline sees only a centroid, while the multi-vector encoding is built from richer handcrafted statistics, including average radial distance, principal orientation planes, and a volume measure. The comparison is suggestive of what an operationally impoverished representation costs, but it does not isolate the effect of discarding the geometric algebra from the effect of discarding the additional summary statistics, and should not be read as a controlled measurement of either alone. Separately, the reported 99\%-encrypted-versus-100\%-plaintext result uses handcrafted rather than gradient-trained weights on a three-class toy problem, a limitation the source paper states outright, with full gradient-based training left as future work. What the case actually supports is the structural claim, that reachable distinctions depend on preserved operations and not merely on preserved values, not the specific accuracy figures generalizing beyond this proof of concept.

\section{Limits and Scope}

The source paper is explicit about where the construction stops paying off. Geometric algebras of dimension $n$ have $2^n$ components, so $\mathrm{Cl}(4,0)$ with 16 components remains feasible, $\mathrm{Cl}(5,0)$ with 32 components is described as borderline, and $\mathrm{Cl}(6,0)$ with 64 components is described as problematic, since ciphertext count grows exponentially with dimension. Problems without underlying geometric structure see no benefit from the construction, and the dense multi-vector representation is reported as inefficient for sparse-matrix problems. Performance remains far from real-time: a single homomorphic multiplication takes approximately 217 milliseconds on the reported CPU implementation against roughly 50 nanoseconds for the equivalent unencrypted geometric product, a slowdown on the order of seven million times, a gap the source paper attributes to the expense of the noise-management operations, relinearization and rescaling, that FHE schemes require after each multiplication.

\section{Conclusion}

Read against the apparatus of Chapter~\ref{ch:ch15-compression-fallacy}, Clifford FHE is not primarily a compression result at all, despite its headline expansion figure inviting that reading. It is a demonstration that a representation's operational repertoire, its native computability over encoded values, is a genuinely independent axis from both how many bits the encoding occupies and how many distinctions it statically preserves, and that this axis can be worth an enormous cost in the other two when what a downstream consumer actually needs is not the values themselves but the structure relating them. This closes Part~IV's argument on the same note it opened on: a representation's apparent quality, whether measured by brevity, by fluency, or by cryptographic compactness, is never a single number, and the cases this Part has worked through, Forth's word libraries, borrowed intuition's provenance gap, SRHS's distributed secret fields, and now Clifford FHE's inverted compression tradeoff, are four instances of the same warning, that the visible artifact rarely reports which of its several independent coordinates was actually optimized, or at what cost to the others.
